\subsection{INNER PRODUCTS: CASE OF BIGEBRAS}

Let E be a graded bigebra (resp. skew graded bigebra) (no. 4, Definition 3); then the results of nos. 6 and 7 can be applied to define the right (resp. left) inner products \(x \llcorner u \in \mathbf{E}\) and \(u \llcorner x \in \mathbf{E}^{*\mathbf{gr}}\) (resp. \(u \lrcorner x \in \mathbf{E}\) and \(x \lrcorner u \in \mathbf{E}^{*\mathbf{gr}}\) )

for all \(x \in \mathbf{E}\) and all \(u \in \mathbf{E}^{*\mathbf{gr}}\) . Thus an (E, E)-bimodule structure and an \((\mathbf{E}^{*\mathbf{gr}}, \mathbf{E}^{*\mathbf{gr}})\) -bimodule structure are obtained on E. Further:

PROPOSITION 9. Let E be a graded bigebra (resp. skew graded bigebra). For every element x of degree 1 in E, the left and right inner products by x are derivations (resp. antiderivations) (§ 10, no. 2) of the algebra E*gr.

In the notation of no. 6, for every homogeneous element \(x\) of degree 1 in a graded bigebra (resp. a skew graded bigebra) E,

\[
c (x) = x \otimes 1 + 1 \otimes x,
\]

by Proposition 5 of no. 3 and the fact that \(c\) is a homomorphism of degree 0. Suppose first that \(E\) is a graded bigebra. For all \(y \in E\) , by definition

\[
\langle (u v) \llcorner x, y \rangle = \langle u v, x y \rangle = m ((u \otimes v) (c (x y)))
\]

and since \(c\) is an algebra homomorphism, \(c(xy) = c(x)c(y)\) . Let \(c(y) = \sum_{i} s_i \otimes t_i\) with \(s_i\) and \(t_i\) in \(E\) ; therefore

\[
c (x y) = \sum_ {i} (x s _ {i}) \otimes t _ {i} + \sum_ {i} s _ {i} \otimes (x t _ {i}).
\]

Hence \(\langle (uv) \llcorner x, y \rangle = \sum_{i} u(xs_i)v(t_i) + \sum_{i} u(s_i)v(xt_i)\) . But we may write

\[
\sum_ {i} u (x s _ {i}) v (t _ {i}) = m (((u \llcorner x) \otimes v) (c (y))) = \langle (u \llcorner x) v, y \rangle
\]

and similarly

\[
\sum_ {i} u (s _ {i}) v (x t _ {i}) = m ((u \otimes (v \llcorner x)) (c (y))) = \langle u (v \llcorner x), y \rangle ,
\]

whence, returning to the notation \(i(x)\) for the inner product,

\[
(i (x)) (u v) = ((i (x)) (u)) v + u ((i (x)) (v))\tag{64}
\]

which proves that \(i(x)\) is a derivation in \(\mathbf{E}^{*}\mathbf{gr}\) .

Suppose now that \(\mathbf{E}\) is a skew graded bigebra, that \(u\in \mathbf{E}_p^*\) , \(v\in \mathbf{E}_q^*\) and \(y\in \mathbf{E}_r\) ; then we can write

\[
c (y) = \sum_ {0 \leqslant j \leqslant r} \left(\sum_ {i} s _ {i j} \otimes t _ {i, r - j}\right)
\]

where the \(s_{ij}\) and \(t_{ij}\) belong to \(\mathbf{E}_j\) ; by definition of the product in \(\mathbf{E}^{\mathfrak{g}} \otimes_{\mathbf{A}} \mathbf{E}\) , then

\[
c (x y) = c (x) c (y) = \sum_ {0 \leqslant j \leqslant r} \left(\sum_ {i} \left(x s _ {i j}\right) \otimes t _ {i, r - j} + (- 1) ^ {j} \sum_ {i} s _ {i j} \otimes \left(x t _ {i, r - j}\right)\right)
\]

whence this time

\[
\langle (u v) \llcorner x, y \rangle = \sum_ {0 \leqslant j \leqslant r} \left(\sum_ {i} u \left(x s _ {i j}\right) v \left(t _ {i, r - j}\right) + (- 1) ^ {j} \sum_ {i} u \left(s _ {i j}\right) v \left(x t _ {i, r - j}\right)\right).
\]

Then also \(\sum_{0\leqslant j\leqslant r}\left(\sum_{i}u(xs_{ij})v(t_{i,r - j})\right) = \langle (u\llcorner x)v,y\rangle\) . On the other hand, \(u(s_{ij}) = 0\) unless \(j = -p\) and hence we may also write

\[
\sum_ {0 \leqslant j \leqslant r} (- 1) ^ {j} \left(\sum_ {i} u (s _ {i j}) v (x t _ {i, r - j})\right) = (- 1) ^ {p} \langle u (v \llcorner x), y \rangle .
\]

We therefore conclude that

\[
(i (x)) (u v) = ((i (x)) (u)) v + (- 1) ^ {p} u ((i (x)) (v)),\tag{65}
\]

in other words \(i(x)\) is an antiderivation in \(\mathrm{E}^{*\mathrm{gr}}\) . The assertions relating to the left inner product by an element \(x\) of degree 1 in E are proved similarly.

Remarks. (1) Let E be a graded bigebra over A and N, N', N'' three graded A-modules. Let m be an A-bilinear mapping of N × N' into N''; for u \ensuremath{\in} Homgr\_A(E, N) and v \ensuremath{\in} Homgr\_A(E, N'), let u.v denote the graded homomorphism m ∘ (u ⊗ v) ∘ c of E into N''. On the other hand, let i(x) denote the (right or left) inner product by x \ensuremath{\in} E in the A-modules Homgr\_A(E, N), Homgr\_A(E, N') and Homgr\_A(E, N'\,'). Then, if x is of degree 1,

\[
(i (x)) (u. v) = ((i (x)) (u)). v + u. ((i (x)) (v))
\]

for all \(u \in \operatorname{Homgr}_{\mathbf{A}}(\mathbf{E}, \mathbf{N})\) and \(v \in \operatorname{Homgr}_{\mathbf{A}}(\mathbf{E}, \mathbf{N}')\) .

Under the same conditions, if E is a skew graded bigebra and u is homogeneous of degree p, then

\[
(i (x)) (u. v) = ((i (x)) (u)). v + (- 1) ^ {p} u. ((i (x)) (v)).
\]

The proofs are the same as in Proposition 9.

\begin{enumerate}
\def\labelenumi{(\arabic{enumi})}
\setcounter{enumi}{1}
\tightlist
\item
  The same argument as in the above proof proves, more generally, that for all \(x \in \mathbf{E}\) , if \(c(x) = \sum_{j} x_j' \otimes x_{j'}'\) then, for all \(u, v\) in \(\mathbf{E}^{*\mathrm{gr}}\) , the ``Leibniz formula''
\end{enumerate}

\[
(i (x)) (u v) = \sum_ {j} (i (x _ {j} ^ {\prime})) (u) \cdot (i (x _ {j} ^ {\prime \prime})) (v)
\]

holds. In particular, for every primitive element of a graded bigebra E, that is such that \(c(x) = x \otimes 1 + 1 \otimes x\) , \(i(x)\) is a derivation of \(\mathrm{E}^{*\mathrm{gr}}\) .

PROPOSITION 10. Let E be a graded bigebra (resp. skew graded bigebra). For every element f of degree 1 in E*gr, the left and right inner products are derivations (resp. antiderivations) of the algebra E.

Let \(x \in \mathrm{E}_p, y \in \mathrm{E}_p\) ( \(p \geqslant 1, q \geqslant 1\) ). By Proposition 5 of no. 3, we can write

\[
c (x) = x \otimes 1 + \sum_ {1 \leqslant j \leqslant p - 1} \left(\sum_ {i} x _ {i j} ^ {\prime} \otimes x _ {i j, p - j} ^ {\prime \prime}\right) + 1 \otimes x
\]

\[
c (y) = y \otimes 1 + \sum_ {1 \leqslant k \leqslant q - 1} \left(\sum_ {i} y _ {i k} ^ {\prime} \otimes y _ {i, q - k} ^ {\prime \prime}\right) + 1 \otimes y
\]

where \(x_{ij}'\) and \(x_{ij}''\) belong to \(\mathrm{E}_j\) , \(y_{ik}'\) and \(y_{ik}''\) to \(\mathrm{E}_k\) . If \(\mathrm{E}\) is a graded bigebra, the component of \(c(xy) = c(x)c(y)\) belonging to \(\mathrm{E}_1 \otimes \mathrm{E}\) , is equal to

\[
\sum_ {i} x _ {i, 1} ^ {\prime} \otimes x _ {i, p - 1} ^ {\prime \prime} y + \sum_ {i} y _ {i, 1} ^ {\prime} \otimes x y _ {i, q - 1} ^ {\prime \prime}
\]

and hence by definition

\[
\begin{array}{r l} (x y) \llcorner f & = \sum_ {i} f (x _ {i, 1} ^ {\prime}) x _ {i, p - 1} ^ {\prime \prime} y + \sum_ {i} f (y _ {i, 1} ^ {\prime}) x y _ {i, q - 1} ^ {\prime \prime} \\ & = (x \llcorner f) y + x (y \llcorner f) \end{array}
\]

and the right inner product by \(f\) is a derivation. If on the other hand \(E\) is a skew graded bigebra, the component of \(c(xy)\) belonging to \(E_1 \otimes E\) is equal to

\[
\sum_ {i} x _ {i, 1} ^ {\prime} \otimes x _ {i, p - 1} ^ {\prime \prime} y + (- 1) ^ {p} \sum_ {i} y _ {i, 1} ^ {\prime} \otimes x y _ {i, q - 1} ^ {\prime \prime}
\]

and this time we obtain

\[
(x y) \llcorner f = (x \llcorner f) y + (- 1) ^ {p} x (y \llcorner f)
\]

which shows that \(i(f)\) is then an antiderivation. The argument is similar for the left inner product by \(f\) .

Examples. Propositions 9 and 10 apply in particular to the graded bigebra \(\mathsf{S}(\mathbf{M})\) and the skew graded bigebra \(\bigwedge (\mathbf{M})\) . The inner products by elements of degree 1 in \(\mathsf{S}(\mathbf{M})\) (resp. \(\mathsf{S}(\mathbf{M})^{*\mathrm{gr}}\) ) are derivations which commute with one another, since \(\mathsf{S}(\mathbf{M})\) (resp. \(\mathsf{S}(\mathbf{M})^{*\mathrm{gr}}\) ) is commutative.

Similarly, the inner products by elements of degree 1 in \(\bigwedge (\mathbf{M})\) (resp. \(\bigwedge (\mathbf{M})^{*\mathrm{gr}}\) ) are antiderivations, which are of zero square, for the square of an element of degree 1 in the algebra \(\bigwedge (\mathbf{M})\) (resp. \(\bigwedge (\mathbf{M})^{*\mathrm{gr}}\) ) is zero.

